\documentclass{beamer}
\usetheme{Madrid}

\usepackage[utf8]{inputenc} % solo si usas pdflatex
\usepackage[T1]{fontenc}    % solo si usas pdflatex

\usepackage{amsmath,amssymb}
\usepackage{microtype}
\usepackage{cancel}         % si lo usas
\usepackage{fontawesome}    % si lo usas
% \usepackage{xypic}        % solo si lo usas
% \usepackage{multicol}     % normalmente no; usar columns de beamer
% \usepackage{amsthm}       % solo si defines teoremas propios
% \usepackage{scrextend}    % solo si lo usas

\usepackage{graphicx}
\graphicspath{{images/}}

\hypersetup{pdfborder={0 0 0}}


\title{Introducción a Kubernetes}
\author{Rodrigo Blanco Betes. Alfa Formación}
\date{Enero del 2026}

\begin{document}

\titlepage

\tableofcontents

\begin{frame}{Introducción a Kubernetes}
\begin{itemize}
\item Kubernetes es una plataforma de orquestación de contenedores de código abierto.
\item Permite automatizar el despliegue, escalado y gestión de aplicaciones containerizadas.
\item Originalmente desarrollado por Google y actualmente mantenido por la Cloud Native Computing Foundation (CNCF).
\item Está diseñado para gestionar aplicaciones distribuidas en clusters de máquinas.
\item Se integra habitualmente con tecnologías como Docker o containerd.
\end{itemize}
\end{frame}

\begin{frame}{Concepto de Cluster}
\begin{itemize}
\item Un cluster de Kubernetes está compuesto por múltiples máquinas llamadas \textbf{nodos}.
\item Los nodos se dividen en dos tipos principales:
\begin{itemize}
\item \textbf{Control Plane}: responsable de la gestión del cluster.
\item \textbf{Worker Nodes}: ejecutan las aplicaciones containerizadas.
\end{itemize}
\item El objetivo del cluster es mantener el estado deseado del sistema.
\item Kubernetes supervisa continuamente el estado real y realiza acciones para corregir desviaciones.
\end{itemize}
\end{frame}

\begin{frame}{Arquitectura general}
\begin{itemize}
\item Kubernetes sigue una arquitectura maestro-trabajador.
\item El \textbf{Control Plane} toma decisiones globales del cluster.
\item Los \textbf{Worker Nodes} ejecutan los contenedores.
\item La comunicación entre componentes se realiza mediante la API de Kubernetes.
\item Todos los recursos del sistema se describen mediante objetos declarativos.
\end{itemize}
\end{frame}

\begin{frame}{Componentes del Control Plane}
\begin{itemize}
\item \textbf{kube-apiserver}
\begin{itemize}
\item Punto de entrada al cluster.
\item Expone la API REST de Kubernetes.
\end{itemize}

\item \textbf{etcd}
\begin{itemize}
\item Base de datos distribuida clave-valor.
\item Almacena el estado completo del cluster.
\end{itemize}

\item \textbf{kube-scheduler}
\begin{itemize}
\item Decide en qué nodo se ejecutará cada Pod.
\end{itemize}

\item \textbf{kube-controller-manager}
\begin{itemize}
\item Ejecuta controladores que garantizan el estado deseado del sistema.
\end{itemize}
\end{itemize}
\end{frame}

\begin{frame}{API Server}
\begin{itemize}
\item El \textbf{kube-apiserver} es el componente central del control plane.
\item Todos los componentes interactúan con el cluster a través de la API.
\item Gestiona operaciones como:
\begin{itemize}
\item creación de recursos
\item modificación de configuraciones
\item consulta del estado del cluster
\end{itemize}
\item Autentica y autoriza las peticiones de los usuarios.
\item También valida los objetos enviados al cluster.
\end{itemize}
\end{frame}

\begin{frame}{etcd}
\begin{itemize}
\item \textbf{etcd} es el almacenamiento persistente de Kubernetes.
\item Guarda toda la configuración del cluster.
\item Mantiene información como:
\begin{itemize}
\item pods
\item servicios
\item deployments
\item configuraciones de red
\end{itemize}
\item Es un sistema distribuido altamente consistente.
\item Su disponibilidad es crítica para el funcionamiento del cluster.
\end{itemize}
\end{frame}

\begin{frame}{Scheduler}
\begin{itemize}
\item El \textbf{kube-scheduler} decide en qué nodo se ejecuta cada Pod.
\item Evalúa múltiples factores:
\begin{itemize}
\item recursos disponibles
\item restricciones de afinidad
\item políticas de scheduling
\item tolerancias y taints
\end{itemize}
\item Selecciona el nodo más adecuado para ejecutar la carga de trabajo.
\end{itemize}
\end{frame}

\begin{frame}{Controller Manager}
\begin{itemize}
\item El \textbf{kube-controller-manager} ejecuta distintos controladores.
\item Los controladores monitorizan el estado del sistema y realizan acciones correctivas.
\item Ejemplos:
\begin{itemize}
\item Node Controller
\item Replication Controller
\item Endpoint Controller
\item Service Account Controller
\end{itemize}
\item Su función principal es mantener el estado deseado del cluster.
\end{itemize}
\end{frame}

\begin{frame}{Componentes de los Worker Nodes}
\begin{itemize}
\item Los worker nodes ejecutan las aplicaciones containerizadas.
\item Componentes principales:
\begin{itemize}
\item \textbf{kubelet}
\item \textbf{kube-proxy}
\item \textbf{container runtime}
\end{itemize}
\item Estos componentes permiten ejecutar, gestionar y conectar los contenedores.
\end{itemize}
\end{frame}

\begin{frame}{Kubelet}
\begin{itemize}
\item El \textbf{kubelet} es el agente que se ejecuta en cada nodo.
\item Se comunica con el API Server.
\item Sus responsabilidades incluyen:
\begin{itemize}
\item ejecutar contenedores
\item monitorizar pods
\item reportar el estado del nodo
\end{itemize}
\item Garantiza que los contenedores definidos en los Pods estén funcionando correctamente.
\end{itemize}
\end{frame}

\begin{frame}{Container Runtime}
\begin{itemize}
\item El container runtime es el software que ejecuta los contenedores.
\item Kubernetes utiliza la interfaz CRI (Container Runtime Interface).
\item Algunos runtimes comunes:
\begin{itemize}
\item containerd
\item CRI-O
\item Docker (históricamente utilizado)
\end{itemize}
\item Se encarga de descargar imágenes, iniciar contenedores y gestionar su ciclo de vida.
\end{itemize}
\end{frame}

\begin{frame}{Kube Proxy}
\begin{itemize}
\item El \textbf{kube-proxy} gestiona la red dentro del cluster.
\item Permite que los servicios sean accesibles desde otros pods.
\item Implementa reglas de red mediante iptables o IPVS.
\item Proporciona balanceo de carga entre pods de un mismo servicio.
\end{itemize}
\end{frame}

\begin{frame}{Pods}
\begin{itemize}
\item El \textbf{Pod} es la unidad mínima de despliegue en Kubernetes.
\item Puede contener uno o varios contenedores.
\item Los contenedores dentro de un pod:
\begin{itemize}
\item comparten red
\item comparten almacenamiento
\item se ejecutan en el mismo nodo
\end{itemize}
\item Los pods son efímeros y pueden ser recreados por Kubernetes.
\end{itemize}
\end{frame}

\begin{frame}{Deployments}
\begin{itemize}
\item Un \textbf{Deployment} describe cómo desplegar y actualizar una aplicación.
\item Permite definir:
\begin{itemize}
\item número de réplicas
\item estrategia de actualización
\item versión de la imagen
\end{itemize}
\item Kubernetes mantiene automáticamente el número deseado de pods.
\item Permite realizar actualizaciones rolling sin interrupciones.
\end{itemize}
\end{frame}

\begin{frame}{Services}
\begin{itemize}
\item Los \textbf{Services} permiten exponer aplicaciones dentro o fuera del cluster.
\item Proporcionan una dirección IP estable para acceder a los pods.
\item Tipos principales:
\begin{itemize}
\item ClusterIP
\item NodePort
\item LoadBalancer
\end{itemize}
\item Permiten balanceo de carga entre múltiples pods.
\end{itemize}
\end{frame}

\begin{frame}{Resumen}
\begin{itemize}
\item Kubernetes permite gestionar aplicaciones distribuidas a gran escala.
\item Su arquitectura se basa en un control plane y múltiples worker nodes.
\item Automatiza tareas clave como despliegue, escalado y recuperación.
\item Es una tecnología fundamental en arquitecturas cloud-native.
\end{itemize}
\end{frame}

\begin{frame}{Networking en Kubernetes}
\begin{itemize}
\item Kubernetes utiliza un modelo de red plano en el que todos los Pods pueden comunicarse entre sí.
\item Cada Pod recibe una dirección IP única dentro del cluster.
\item La comunicación entre Pods no requiere NAT.
\item El networking se implementa mediante plugins llamados \textbf{CNI (Container Network Interface)}.
\item Algunos plugins CNI populares incluyen:
\begin{itemize}
\item Calico
\item Flannel
\item Cilium
\item Weave Net
\end{itemize}
\end{itemize}
\end{frame}

\begin{frame}{Servicios y exposición de aplicaciones}
\begin{itemize}
\item Los \textbf{Services} permiten exponer pods y proporcionar una dirección estable.
\item Kubernetes implementa balanceo de carga entre las réplicas de pods.
\item Tipos principales de servicios:
\begin{itemize}
\item \textbf{ClusterIP}: accesible solo dentro del cluster.
\item \textbf{NodePort}: expone el servicio en un puerto del nodo.
\item \textbf{LoadBalancer}: utiliza un balanceador de carga externo.
\end{itemize}
\item Los \textbf{Ingress Controllers} permiten gestionar acceso HTTP/HTTPS con reglas de enrutamiento.
\end{itemize}
\end{frame}

\begin{frame}{Almacenamiento en Kubernetes}
\begin{itemize}
\item Los contenedores son efímeros, por lo que Kubernetes proporciona mecanismos de almacenamiento persistente.
\item Los principales recursos son:
\begin{itemize}
\item \textbf{Persistent Volumes (PV)}
\item \textbf{Persistent Volume Claims (PVC)}
\item \textbf{Storage Classes}
\end{itemize}
\item Un \textbf{Persistent Volume} representa almacenamiento físico o virtual disponible en el cluster.
\item Un \textbf{Persistent Volume Claim} solicita almacenamiento con determinadas características.
\item Las \textbf{Storage Classes} permiten aprovisionamiento dinámico de almacenamiento.
\end{itemize}
\end{frame}

\begin{frame}{Escalado automático}
\begin{itemize}
\item Kubernetes permite escalar aplicaciones automáticamente según la carga.
\item Tipos principales de autoscaling:
\begin{itemize}
\item \textbf{Horizontal Pod Autoscaler (HPA)}
\begin{itemize}
\item Ajusta el número de pods según métricas como CPU o memoria.
\end{itemize}
\item \textbf{Vertical Pod Autoscaler (VPA)}
\begin{itemize}
\item Ajusta los recursos asignados a un pod.
\end{itemize}
\item \textbf{Cluster Autoscaler}
\begin{itemize}
\item Añade o elimina nodos del cluster automáticamente.
\end{itemize}
\end{itemize}
\end{itemize}
\end{frame}

\begin{frame}{Seguridad en Kubernetes}
\begin{itemize}
\item Kubernetes incluye varios mecanismos para proteger el cluster.
\item \textbf{RBAC (Role Based Access Control)} permite definir permisos de usuarios y servicios.
\item Los \textbf{Namespaces} permiten aislar recursos dentro del cluster.
\item Los \textbf{Network Policies} controlan el tráfico entre pods.
\item Los \textbf{Secrets} permiten almacenar información sensible como contraseñas o tokens.
\end{itemize}
\end{frame}

\begin{frame}{Gestión declarativa}
\begin{itemize}
\item Kubernetes utiliza un modelo declarativo.
\item Los usuarios describen el estado deseado del sistema mediante archivos YAML.
\item Kubernetes compara continuamente el estado actual con el estado deseado.
\item Si detecta diferencias, realiza acciones para corregirlas.
\item Esto permite automatizar despliegues y facilitar la reproducibilidad.
\end{itemize}
\end{frame}

\begin{frame}{Flujo típico de despliegue}
\begin{enumerate}
\item El desarrollador crea imágenes de contenedor.
\item Las imágenes se almacenan en un registro de contenedores.
\item Se define la aplicación mediante archivos YAML (Deployments, Services, etc.).
\item Se aplican los manifiestos mediante \texttt{kubectl}.
\item Kubernetes programa los pods y gestiona su ejecución.
\item El sistema monitoriza continuamente el estado del cluster.
\end{enumerate}
\end{frame}

\end{document}
